\documentclass[11pt, a4paper]{article}
\usepackage[utf8]{inputenc}
\usepackage[T1]{fontenc}
\usepackage{geometry}
\usepackage{xcolor}
\usepackage{enumitem}
\usepackage{titlesec}
\usepackage{hyperref}

% 页面设置
\geometry{left=2.5cm, right=2.5cm, top=2.5cm, bottom=2.5cm}

% 颜色定义
\definecolor{darkblue}{RGB}{0, 0, 139}
\definecolor{darkgreen}{RGB}{0, 100, 0}
\definecolor{darkred}{RGB}{139, 0, 0}

% 自定义命令
\newcommand{\reviewer}[1]{\vspace{0.5cm}\noindent\textbf{\large \textcolor{darkblue}{#1}}}
\newcommand{\comment}[1]{%
    \par\vspace{1em}% 这里控制空行的大小，可以用 1em, 0.5cm, 10pt 等
    \noindent\textbf{\textit{Comment:}} \textit{\textcolor{gray}{#1}}%
}
\newcommand{\response}[1]{\noindent\textbf{Response:} #1}
\newcommand{\action}[1]{\noindent\textbf{\textcolor{darkgreen}{Action Taken:}} #1}

\title{\textbf{Explanation of Revisions}}
\date{}

\begin{document}

\maketitle

\section{Summary of Major Revisions}
We have carefully revised our manuscript to incorporate the structural suggestions of the Area Chair and to address the constructive feedback from the reviewers.

\begin{description}
\item[1. Restoring the ``Related Work'' Section (Response to Area Chair \& Reviewer BwYb)] \hfill 
\textbf{Location: Section 2} \\
We have moved the ``Related Work'' section from the appendix back to the main body. This ensures the manuscript is self-contained and adheres to standard formatting norms.

\item[2. Expansion of Literature Review (Response to Reviewer BwYb)] \hfill \\
\textbf{Location: Section 2 \& References} \\
We have significantly expanded the bibliography to include foundational, peer-reviewed ASG publications spanning from 2019 to 2022. This ensures a comprehensive coverage of the field's development.

\item[3. Theoretical Mapping (Response to Reviewer BwYb)] \hfill \\
\textbf{Location: Appendix H (Table 13)} \\
We incorporated a comprehensive theoretical mapping table to explicate the alignment between our 16 metrics and the specific cognitive levels of Bloom's Taxonomy. This mapping strengthens the theoretical rationale behind our proposed framework.

\item[4. Validation against Human Experts] \hfill \\
\textbf{Location: Appendix G.1} \\
Addressing concerns regarding metric validity, we added a human-system correlation study. With Pearson coefficients exceeding 0.7, we empirically demonstrate that our metrics align with human evaluations.

\item[5. Diagnostic Analysis (Response to Reviewer zYdN)] \hfill \\
\textbf{Location: Appendix F} \\
We introduced a ``Diagnostic Analysis of Performance Bottlenecks'' section to analyze the underlying causes of the observed performance. We explicitly attribute failures to systemic issues, including retrieval blindness, lossy compression, and meta-knowledge deficit.

\end{description}
\section{Response to the Area Chair}

\reviewer{Response to Area Chair}

\comment{``It is not acceptable to move related work discussion into an Appendix!! ... Authors should resubmit a revised version of the paper where the main paper body includes related work and their response to above issues and concerns. If they do this well, this will be a strong paper which deserves to be presented at xACL.''}

We thank the Area Chair for the encouraging assessment and the helpful guidance on structural norms. We have fully implemented the directive. To address the concern regarding the feasibility of fitting all revisions within the page limit, we adopted a strategic content reorganization:

\begin{itemize}
    \item \textbf{Strict Adherence to Structure (Section 2):} We have restored the \textbf{Related Work} section to the main body (Section 2), ensuring the paper is self-contained.
    
    \item \textbf{Balancing Depth and Space:} To accommodate both the Related Work and the comprehensive responses to reviewers (BwYb \& zYdN) without exceeding the page limit, we prioritized core insights in the main text while placing supporting evidence in the appendices.
    Specifically:
    \begin{itemize}
        \item \textbf{Main Text:} Now includes the expanded bibliography (Section 2) and synthesis of key results.
        \item \textbf{Appendices (Referenced in Text):} We utilized the appendix for extensive data that supports our claims, including the human validation study (Appendix G.1), theoretical mapping (Appendix H), and diagnostic root cause analysis (Appendix F).
    \end{itemize}

    \item Further details on these specific changes are provided in the Response to Reviewers section below.
\end{itemize}


\section{Response to Reviewers}
\reviewer{Response to Reviewer BwYb}

\comment{There are some problems with the references of this work. About 50\% of them are pre-prints... At the same time, most non pre-print refs are around 10 years old...}

We thank the reviewer for the rigorous scrutiny of our bibliography. We have conducted a thorough audit to address this:

\begin{itemize}
    \item \textbf{Bibliography Enrichment:} We have significantly expanded our coverage of peer-reviewed literature, specifically adding relevant works published between 2019 and 2022. This bridges the gap between earlier works and recent LLM-based advancements.
    \item \textbf{Updating and Justifying Pre-prints:} We updated citations for papers recently accepted to venues (e.g., updating MLR-bench \cite{chen2025mlrbench} to its NeurIPS version). While we prioritize peer-reviewed sources, we retained select high-impact pre-prints (e.g., SurveyX \cite{liang2025surveyx}) solely to benchmark against the rapidly evolving state-of-the-art, which is standard practice in the fast-moving LLM field.
\end{itemize}

\comment{Related works are reported in an appendix. It's the first time I see this...}

We fully agree with this assessment. We have moved the Related Work section from the appendix back to the main body (Section 2).

\comment{Will the benchmark be released? There is no link (even anonymised) or mention to release in the paper.}

We are committed to open science. We have added an explicit statement in Section 1 to confirm the availability of resources. The text reads: ``We will open-source the benchmark and evaluation scripts upon publication to foster future research.''

\comment{Lines 066-067: moving from ``formal mimicry'' to ``cognitive deepening''. What does this mean? There's no additional detail.}

We expanded the phrasing in the abstract to explicitly define the contrast. The text now reads: ``...shifting the research focus from the formal mimicry of surface structure to the cognitive deepening of intellectual content.''


\comment{Lines 119-122: Paraphrasing, the LLM is either absent, used for a specific task, or used as a judge... So what's the difference in this approach?}

We clarify that the novelty lies not in the roles themselves, but in the dual-constraint strategy that strictly assigns these roles based on cognitive levels to maximize reliability. We have explicitly articulated this principle in Section 3.1. The revised text reads: ``Our framework is built upon a dual-constraint strategy... This strategy strictly regulates the operational boundary of the LLM based on task subjectivity. For quantifiable metrics, the model is either entirely absent... or confined to specific functional tools... for abstract cognitive dimensions, it operates solely as a constrained judge.''

\comment{Section 2.1.2: Is there a reason why the authors introduce a new metric for Distributional Similarity, instead of using the pre-existing ones?}

We clarify that Distributional Similarity (DS) is not a newly invented mathematical formula, but a composite metric integrating three established distances to ensure robustness against outliers. We have revised the definition in Appendix A.1 to explicitly list the standard components. The text now reads: ``The DS score is ... a multi-metric fusion of three distances empirically validated in NLP contexts: the Jensen-Shannon Divergence, the Hellinger Distance, and the Total Variation Distance.''


\comment{Section 2.1.3. This approach looks like a 2-step LLM-as-judge. What exactly makes it more reliable... And can the authors say that this transforms a black box into something auditable? Also, some examples... And what prompt was used?}

We appreciate this opportunity to clarify the core contribution of our GRADE approach. We have addressed the following sub-questions in the revised manuscript:

\begin{itemize}
    \item \textbf{Reliability \& Auditability (Section 3.1):} We clarified that GRADE differs from standard judges by forcing the model to externalize its criteria (Step 1: Rubric Generation) before scoring. This transforms an unexplainable scalar into a transparent, diagnosable argument, allowing humans to inspect exactly why a score was given.
    \item \textbf{Qualitative Examples (Appendix G.2.3):} To demonstrate this auditability, we added a case study in Appendix G.2.3. It visually contrasts the black-box output of baseline judges with the detailed, evidence-based rubric generated by GRADE.
    \item \textbf{Prompts (Appendix I):} We confirmed that we performed extensive, task-specific prompt engineering tailored to Bloom's Taxonomy rather than reusing generic prompts. The full set of prompts is provided in Appendix I.
\end{itemize}

\comment{It is unclear to me how Bloom-Eval reflects the Bloom taxonomy... The only thing I can see is that the ``levels'' of the evaluation resemble it in their naming.}

We added a theoretical mapping table in Appendix H (Table 13) that systematically aligns each of our 16 metrics with the specific cognitive definitions of Bloom's Taxonomy, justifying their diagnostic purpose.

\comment{Regarding HIRC: doesn't this metric bias the model towards retrieving only highly-cited work, hence potentially excluding recently released relevant papers?}

We acknowledge the valid point that HIRC favors classic works. We consider that a survey should fundamentally cover classics. The coverage of recent papers is implicitly captured by our Entity Fidelity (EFid) metric. Since our benchmark targets 2023-2025 papers, missing recent papers results in missing key new entities, which is reflected in the EFid score.

\comment{FCons: So, this is an LLM being trusted to semantically match the claims? How is this reliable? Again, prompt? How was it validated? The same applies to OTC.}

To establish reliability, we benchmarked these metrics against human expert ratings in our new correlation study. The results (Appendix G.1) show a strong alignment (Pearson correlation of 0.76 for Factual Consistency), demonstrating that the LLM's semantic matching serves as a reliable proxy for human judgment. We have released the full, specific prompts used for both FCons and OTC in Appendix I, allowing for complete reproducibility and scrutiny of the matching logic.

\comment{Across all levels, how are all these metrics justified? ... And how are they combined? ... This is concerning, especially considering this paper introduces 18 different metrics.}

We agree that such a comprehensive metric suite requires rigorous grounding. In the revised manuscript, we have further strengthened the theoretical justification by incorporating a theoretical mapping table (Table 13), which explicitly aligns each metric with the specific cognitive levels of Bloom's Taxonomy. Furthermore, we clarify that the metrics are combined using an unweighted arithmetic mean to maintain a neutral perspective, as specified in the Figure 4 caption.


\comment{There are no human judgments for most of the metrics; how can we believe the metric actually works? Also, the paper mentions human experts at various steps but it's unknown who these experts are, their background, how they produced the references etc.}

Addressing concerns regarding metric validity, we added a human-system correlation study in Appendix G.1. With Pearson coefficients exceeding 0.7, we empirically demonstrate that our metrics are consistent with human evaluations. We explicitly defined the ``experts'' referenced in the paper. In Section 4.1, we specify that for the dataset ground truth, the experts are the original authors of the 3,506 peer-reviewed survey papers (selected from top-tier venues). In Appendix G.1, we detail that for the human evaluation study, the experts consist of 3 PhD candidates and 2 Master's students.

\comment{What's ``Automatic Survey Generation''? ...}

We have added a formal definition of ASG in Section 1 to address this. We explicitly defined the task as follows: ``ASG is the automated process of synthesizing relevant, disparate literature into a survey given a specific research topic.''

\comment{Lines 127-136: I don't think that in 2025 an entire section is needed to define the F1-score. It's an extremely well-known measure that can just be mentioned.}

We have removed the explicit mathematical definition of the F1-score to optimize space, simply mentioning it as a standard metric.

\comment{Lines 138-142: paraphrasing, ``to measure distributional similarity, we propose a metric called distributional similarity''. I think this can be worded a bit better.}

We revised the description in Appendix A.1 to read: ``To quantify the distributional alignment between sets (e.g., entities or topics), we propose a composite metric named Distributional Similarity (DS).''

\comment{Figure 3: Can you please correct the scale? Your highest value barely touches half of it, and as a result, the other benchmarks you're comparing against are invisible.}

We have redrawn Figure 3 with an optimized scale.

\reviewer{Response to Reviewer 5wHA}

\comment{Inherent Subjectivity: The GRADE method, while transparent, still relies on an LLM, which has ``inherent subjective preferences'' and is subject to ``Rater Bias.''}

Our human-system correlation study (Appendix G.1) shows strong alignment ($>0.7$) with human experts. This empirically demonstrates that, despite inherent model subjectivity, the resulting evaluation serves as a reliable proxy for human consensus. Furthermore, the transparent nature of GRADE allows for human oversight, mitigating the risks associated with opaque rater bias.

\comment{Limited Benchmark Scope: The large benchmark may still not fully represent all scientific research domains, limiting the generalizability of the findings.}

While our benchmark covers over 3,000 papers across Computer Science and 51 high-impact journal series, we agree it cannot cover every niche domain. However, we intentionally prioritized quality and depth (peer-reviewed, high-impact venues) over sheer quantity (e.g., unverified arXiv preprints) to ensure the evaluation standard is authoritative and representative of top-tier academic work. 

We view Bloom-Eval as a continuously evolving project and are committed to expanding our data sources to include a broader range of disciplines in future iterations.

\reviewer{Response to Reviewer zYdN}

\comment{The paper's own sensitivity analysis (Appendix E) shows only 10\% directional agreement between judges and systematic bias patterns (optimistic/pessimistic/unstable raters).}

We clarify that this sensitivity analysis was intended to visualize the inherent challenge and subjectivity of assessing high-order cognitive skills rather than a failure of our methodology. To resolve the concern regarding reliability, we have added a human-system correlation study in Appendix G.1. The results demonstrate that our GRADE approach achieves strong alignment (Pearson $>0.7$) with human experts, proving that it serves as a reliable proxy for human judgment despite the systematic biases of individual LLM raters. Furthermore, our findings suggest that the selection of an appropriate Judge LLM is crucial. Practitioners should prioritize models that demonstrate high consistency with human expertise (as we identified with GPT-5-mini) to ensure that the evaluation remains authoritative.

\comment{The paper extensively documents failure patterns but provides no analysis of root causes ... and the equal-weight metric aggregation ignores which cognitive dimensions matter most for practical deployment.}


We have addressed this by adding a detailed diagnostic analysis of performance bottlenecks in Appendix F. In this section, we analyze the root causes of failure, such as retrieval misalignment and lossy information compression in current ASG architectures. Furthermore, the inherent hierarchical structure of Bloom-Eval is specifically designed to provide fine-grained diagnostics. While we report the unweighted arithmetic mean as a neutral baseline for overall benchmarking, the modularity of our framework empowers practitioners to easily assign custom weights to different cognitive dimensions (e.g., prioritizing Evaluation over Memory) to align with their specific deployment requirements.

\reviewer{Response to Reviewer 5HGw}

We gratefully acknowledge the Area Chair's decision to exclude this review from consideration due to factual inaccuracies.

\bibliography{custom}
\bibliographystyle{unsrt}
\end{document}